\vsssub
\subsubsection{~$S_{in} + S_{ds}$：Rogers et al. 2012 \& Zieger et al. 2015} \label{sec:ST6}
\vsssub

\opthead{ST6}{AUSWEX, Lake George}{A. Babanin, I. Young, M. Donelan, E. Rogers, S. Zieger, Q. Liu}

\noindent
这一版本实现了基于观测的深水源汇项物理过程。这些过程包括风输入源项，以及由负风输入、
白帽耗散和波-湍流相互作用（涌浪耗散）造成的汇项。风输入和白帽耗散源项基于澳大利亚
Lake George 的观测；波-湍流耗散基于实验室实验和涌浪衰减的现场观测；负输入基于
实验室测试。总风能输入受到风应力约束，而风应力可独立获得。

\paragraph{风输入。} 除了首次直接现场测量强风强迫下的风输入外，Lake George 实验还揭示了
此前未被考虑的若干风-浪交换新物理特征：
(i) 完全气流分离，它会在强风/陡波条件下导致风输入相对减小；
(ii) 波增长率对波陡的依赖，这表明风输入源函数具有非线性行为；
(iii) 存在破碎波时输入增强
\citep{art:Dea06,art:Bea07}（最后这一特征未在这里实现）。
按照 \citet{art:RBW12}，该输入源项写为
\begin{eqnarray}
\cS_{in}(k,\theta) & = & \frac{\rho_a}{\rho_w}\, \sigma\,\gamma(k,\theta)\,N(k,\theta)  ,
\label{eq:ST601} \\ \gamma(k,\theta) &=& G\,\sqrt{B_n}\,W  ,
\label{eq:ST602} \\
G                &=& 2.8-\Bigl (1+\tanh (10\sqrt{B_n}W-11) \Bigr)  ,
\label{eq:ST603} \\
B_n              &=& A(k)\,N(k)\sigma\,k^3  ,
\label{eq:ST604} \\
W                &=& \left (\frac{U_s}{c}\,-\,1 \right )^2  .
\label{eq:ST605}
\end{eqnarray}

\noindent
在式 (\ref{eq:ST601})--(\ref{eq:ST605}) 中，$\rho_a$ 和 $\rho_w$ 分别为空气和水的密度，
$U_s$ 为尺度风速，$c$ 表示波相速，$\sigma$ 为弧频率，$k$ 为波数。
\citet{art:Phi84} 引入的谱饱和度 (\ref{eq:ST604}) 是波陡 $ak$ 的谱度量。
全方向作用量密度通过对所有方向积分得到：
$N(k)=\int N(k,\theta)d\theta$。\linebreak
方向谱狭窄度的倒数 $A(k)$ 定义为\linebreak
$A^{-1}(k) =$ $\int_{0}^{2\pi} [{N(k,\theta)}/{N_{\max}(k)}] d\theta$，
其中 $N_{\max}(k)=\max\bigl \{N(k,\theta)\bigr \}$，对所有方向
$\theta\in[0,2\pi]$ 取最大值 \citep{art:BS87}。

\citet{art:Dea06} 用平均海面以上 10\,m 风速来参数化增长率 (\ref{eq:ST602})。
然而，波浪模型通常使用摩擦速度 $u_\star=\tau/\rho_a$，以保证不同风速下风区律的一致性
\citep[][p. 253]{bk:WAM94}。因此，\citet{art:RBW12} 按照 \citet{art:KHH84}，
主张使用近似
\begin{equation}
U_s = U_{10} \simeq \Upsilon u_{\ast},\ \mathrm{and}\ \Upsilon = 28
\label{eq:Upsilon}
\end{equation}

\begin{eqnarray}
W_1 & = & \mathrm{max}^2 \left \{ 0,\frac{U}{c}\ \cos(\theta-\theta_w)-1
\right \}  , \label{eq:ST606-1} \\
W_2 & = & \mathrm{min}^2 \left \{ 0,\frac{U}{c}\ \cos(\theta-\theta_w)-1
\right \} .\label{eq:ST606-2}
\end{eqnarray}

\noindent
$W$ 的方向分布实现为顺风项 (\ref{eq:ST606-1}) 和逆风项 (\ref{eq:ST606-2}) 之和，
使二者互补（即 $W=\{W_1\cup W_2$\}；见本节后文的{\it 负输入}）：
\begin{equation}\label{eq:ST606}
W=W_1-a_0\,W_2  .
\end{equation}

\paragraph{风输入约束。} 输入项的一个重要部分是计算从大气到海洋的动量通量，
它必须与波接收的通量一致。在海面上，应力 $\vec{\tau}$ 可写为黏性应力和波支撑应力之和：
$\vec{\tau} = \vec{\tau}_{v} + \vec{\tau}_{w}$。波支撑应力 $ \vec{\tau}_{w}$
用作风输入的主要约束，且不能超过总应力 $\vec{\tau} \le \vec{\tau}_{tot}$。
这里总应力由通量参数化确定：$\vec{\tau}_{tot}=\rho_a u_\star|u_\star|$。
波支撑应力 $\tau_w$ 可通过对风动量输入函数积分来计算：

\begin{equation}\label{eq:ST609}
   \vec{\tau}_w = \rho_w g \int_{0}^{2\pi} \int_{0}^{k_{max}}
   \frac{S_{in}(k^\prime,\theta)}{c} \Bigl (\cos\theta,\sin\theta \Bigr )
   dk^\prime d\theta  .
\end{equation}

\noindent
波支撑应力 (\ref{eq:ST609}) 的计算包括直到最高离散波数 $k_{max}$ 的已解析谱部分，
也包括短波支撑的应力。为了考虑后者，在能量密度谱最高频率之外假定一条 $f^{-5}$
诊断尾部。为满足约束，在 $\vec{\tau} > \vec{\tau}_{tot}$ 的情况下，施加依赖波数的
因子 $L$ 来削减谱高频部分的能量：$S_{in}(k^\prime)=L(k^\prime)\,S_{in}(k^\prime)$，
其中

\begin{equation}\label{eq:ST610}
L(k^\prime) = \min \Bigl \{ 1, \exp \bigl ( \mu\,[1- U/c] \bigr ) \Bigr
\}  .
\end{equation}

\noindent
削减式 (\ref{eq:ST610}) 是风速和相速的函数，并采用指数形式，旨在削减谱离散部分的能量。
削减强度由系数 $\mu$ 控制，它对高频影响较大，而对能量占优部分影响很小。
$\mu$ 的取值在每个积分时间步通过迭代动态计算 \citep{art:Tea10}。

阻力系数为
\begin{equation}\label{eq:ST607}
C_d \times 10^4 = 8.058 + 0.967 U_{10} - 0.016 U_{10}^2 ,
\end{equation}
该式被选用并实现为开关 {\code FLX4}。这一参数化由 \citet{art:Hwa11} 提出，
考虑了风速超过 30\,m~s$^{-1}$ 时海面阻力的饱和及进一步下降。为防止在极强风
($U_{10}\ge50.33$m~s$^{-1}$) 下 $u_\star$ 降为零，表达式 (\ref{eq:ST607})
被修改为给出 $u_\star=2.026$m~s $^{-1}$。{\it 重要！} 在 {\code ST6} 中，
对风输入场任一均匀偏差的整体调整，是通过风应力参数 $u_\star$ 而不是 $U_{10}$ 完成的。
为此，(\ref{eq:ST607}) 左端表达式 $C_d \times 10^4$ 中的因子被替换为
$C_d \times \mathrm{FAC}$，并作为 {\F FLX4} namelist 参数 {\code CDFAC} 加入
（见本节末尾的{\it 整体调整}）。黏性阻力系数
\begin{equation}\label{eq:ST608}
  C_v \times 10^3 = 1.1 - 0.05 U_{10} ,
\end{equation}
由 \citet{art:Tea10} 使用 \citet{art:BP98} 的数据参数化为风速函数。

\paragraph{负输入。} 除正输入外，{\code ST6} 还包含负输入项，用于在波谱中存在逆向
风应力分量的部分抑制波增长 (\ref{eq:ST606-1})--(\ref{eq:ST606-2})。逆风下的增长率
为负 \citep{pro:Don99}，并在满足波支撑应力 $\tau_w$ 的约束后施加。
(\ref{eq:ST606}) 中 $a_0$ 的取值是输入参数化中的调节参数，可通过 {\F SIN6}
namelist 参数 {\code  SINA0} 调整。

\noindent
\paragraph{白帽耗散。} 对于由波浪破碎造成的耗散，Lake George 现场研究揭示了若干
新特征：(i) 波浪破碎的阈值行为 \citep{art:BBY01}。波浪只有在超过某一通用波陡时
才会破碎，此时破碎概率取决于超过该阈值波陡的程度。低于临界阈值的波，白帽耗散为零。
(ii) 由于长波影响短波破碎和耗散而产生的累积耗散效应
\citep{pro:Don01,pro:BY05,art:Mea06,art:YB06,art:Bea10}；
(iii) 强风下的非线性耗散函数
(\citeauthor{art:Mea06}, \citeyear{art:Mea06};\linebreak
 \citeauthor{art:Bea07}, \citeyear{art:Bea07})；
(iv) 耗散方向展宽的双峰分布
\citep{art:YB06,art:Bea10}（最后这一特征未在 {\code ST6} 中实现）。
按照 \citet{art:RBW12}，白帽耗散项实现为
\begin{equation}\label{eq:ST620}
  \cS_{ds}(k,\theta) = \Bigl [ T_1(k,\theta) + T_2(k,\theta) \Bigr ]\ N(k,\theta) ,
\end{equation}

\noindent
其中 $T_1$ 是固有破碎项，表示为传统的波谱函数；$T_2$ 表示为低于波数 $k$ 的波谱积分，
用于描述每个频率/波数上由长波造成的短波破碎或耗散的累积效应。当该频率处于谱峰或低于
谱峰时，固有破碎项 $T_1$ 是唯一的破碎耗散项。一旦谱峰移动到该频率以下，$T_2$
就会被激活，并随谱峰继续下移而逐渐变得更重要。

阈值谱密度 $F_{\mathrm{T}}$ 计算为
\begin{equation}\label{eq:ST621}
  F_{\mathrm{T}}(k)=\frac{\varepsilon_{\mathrm{T}}}{A(k)\,k^3}  ,
\end{equation}
其中 $k$ 为波数，$\varepsilon_{\mathrm{T}}=0.035^2$ 为经验常数
\citep{art:Bea07,bk:Bab11}。令超过临界阈值谱密度（此时波浪破碎占主导）的超越量为
$\Delta(k)=F(k)-F_{\mathrm{T}}(k)$。进一步令 $\mathcal{F}(k)$ 为用于归一化的通用谱密度，
则固有破碎分量可计算为
\begin{equation}\label{eq:ST622}
T_1(k)=a_1 A(k)\frac{\sigma}{2\pi} \left [ \frac{\Delta(k)}{\mathcal{F}(k)}
\right ]^{p_1}  .
\end{equation}

\noindent
累积耗散项在频率空间中是非局地的，并基于一个向高频增长的积分，在较小尺度上占主导：

\begin{equation}\label{eq:ST623}
T_2(k)=a_2 \int\limits_0^k A(k) \frac{c_g}{2\pi} \left [
\frac{\Delta(k)}{\mathcal{F}(k)} \right ]^{p_2}\!\!dk .
\end{equation}

\noindent
耗散项 (\ref{eq:ST622})--(\ref{eq:ST623}) 依赖五个参数：用于归一化的通用谱密度
$\mathcal{F}(k)$，以及四个系数 $a_1$、$a_2$、$p_1$ 和 $p_2$。系数 $p_1$ 和 $p_2$
控制耗散项中归一化阈值谱密度 $\Delta(k)/\mathcal{F}(k)$ 的强度。Namelist 参数
{\code SDSET} 在 (\ref{eq:ST622})--(\ref{eq:ST623}) 的归一化中切换使用谱密度
$F(k)$ 或阈值谱密度 $F_{\mathrm{T}}(k)$。根据 \citet{art:Bea07} 和
\citet{art:Bab09}，式 (\ref{eq:ST621})--(\ref{eq:ST623}) 中方向狭窄度参数取为
一，即 $A(k)\approx 1$。

\begin{table} \begin{center}
\footnotesize
\begin{threeparttable}
\begin{tabular}{|l|c|c|c|c|c|} \hline \hline
参数              &  WWATCH 变量 & namelist &  vers.\,4.18 & vers.\,5.16 & vers.\,6.07 \\
\hline
  $F_{\mathrm{T}}$ &  SDSET       & SDS6     &  T       &  T         & T          \\
  $a_1$            &  SDSA1       & SDS6     & 6.24E-7  & 3.74E-7    & 4.75E-6    \\
  $p_1$            &  SDSP1       & SDS6     &  4       &  4         & 4          \\
  $a_2$            &  SDSA2       & SDS6     & 8.74E-6  & 5.24E-6    & 7.00E-5    \\
  $p_2$            &  SDSP2       & SDS6     &  4       &  4         & 4          \\
\hline
  $\Upsilon$\tnote{\textdagger}       &  SINWS       & SIN6     & n/a      & n/a        & 32.0       \\
  $N_{hf}$\tnote{\textdagger}         &  SINFC       & SIN6     & n/a      & n/a        & 6.0        \\
  $a_0$            &  SINA0       & SIN6     & 0.04     & 0.09       & 0.09       \\
  $b_1$ 为常数     &  CSTB1       & SWL6     & n/a      &  F         & F          \\
  $b_1$, $B_1$     &  SWLB1       & SWL6     & 0.25E-3  & 0.0032     & 0.0041     \\
  $\mathrm{FAC}$   &  CDFAC       & FLX4     & 1.00E-4  & 1.00E-4    & 1.0\tnote{\textdaggerdbl}        \\
\hline
  $C$              &  NLPROP      & SNL1     & 3.00E7   &  3.00E7    & 3.00E7    \\
 \hline \hline
\end{tabular}
\begin{tablenotes}
	\item[\textdagger] 在 WW3 版本 4.18 和 5.16 中，$\Upsilon = 28.0$
            与 $N_{hf} = 6.0$ 被硬编码在 {\code ST6} 模块中。
    \item[\textdaggerdbl] 对 2011 年及以后的 CFS 风场，全球模拟推荐
            $\mathrm{CDFAC} \sim 1.0$；对 ERA5 再分析风场，建议
            $\mathrm{CDFAC} \sim 1.08$ \citep{Liu2019, Liu2021}。
\end{tablenotes}
\end{threeparttable} \end{center}
\caption{{\code ST6} 同 {\code DIA} 非线性求解器（第~\ref{sec:NL1} 节）联用时的
         率定参数汇总。表中数值代表模型默认设置。缩写 ``n/a'' 表示该变量在对应代码
         版本中不适用。}
\label{tab:ST601} \botline \end{table}

\citet{art:RBW12} 基于时长受限的学术测试率定了耗散项。{\code ST6} 中使用并列于
表~\ref{tab:ST601} 的率定系数，与 \citet{art:RBW12} 的值略有不同，主要原因是
波支撑应力 $\vec{\tau}_w$ 以矢量分量形式实现，并且此前 \citet{art:RBW12} 未考虑
非破碎涌浪耗散 (\ref{eq:ST624})。

\paragraph{涌浪耗散。} 在没有波浪破碎时，还存在其他波浪衰减机制。这里把它们称为
涌浪耗散，并按照波浪同海洋湍流的相互作用来参数化 \citep{bk:Bab11}。不过，该机制
对风生波也保持活跃。如果谱高于波浪破碎阈值，它在整个谱上的贡献较小；但在谱前缘，
甚至在完全 Pierson-Moscowitz 发展情况下的谱峰处，它会占主导。

\begin{equation}\label{eq:ST624}
  \cS_{swl}(k,\theta) = -\frac{2}{3}b_1 \sigma\ \sqrt{B_n}\ N(k,\theta).
\end{equation}

\noindent
通过使式 (\ref{eq:ST624}) 中的系数 $b_1$ 依赖波陡，可以减小波高空间偏差中的大梯度：

\begin{equation}\label{eq:ST625}
   b_1 = B_1 \, 2\sqrt{E}\,k_p .
\end{equation}

在式 (\ref{eq:ST625}) 中，$B_1$ 是尺度系数，$E$ 是式 (\ref{eq:etot}) 中的海面
总方差，$k_p$ 是谱峰波数。式 (\ref{eq:ST625}) 可通过 {\F SWL6} namelist 参数
{\code CSTB1} 启用。式 (\ref{eq:ST625}) 中系数 $B_1$ 和/或式 (\ref{eq:ST624})
中 $b_1$ 的值，可通过 {\F SWL6} namelist 参数 {\code SWLB1} 定制
（见表~\ref{tab:ST601}）。

\noindent
\paragraph{vers.\,6.07 以来的更新} 按照 \citet{art:RBW12}，版本 4.18 和 5.16
采用了尺度风速 $U_s = 28u_{\ast}$ (\ref{eq:Upsilon})（表~\ref{tab:ST601}）。
这些 {\code ST6} 配置已被证明在不同空间尺度和不同天气条件下具有良好技巧
\citep[e.g.][]{art:ZBRY15, liu2017}。然而，\citet[][their Fig. 5]{art:ZBRY15}
也指出 {\code ST6}（版本 4.18 和 5.16）倾向于高估谱高频尾部的能级，表明该特定
频率范围内不同源项之间的平衡不够准确。

Rogers（2014，未发表工作）发现，在 SWAN 的 {\code ST6} 实现中使用
$U_s = 32u_{\ast}$（即 $\Upsilon = 32$）可提高模型估计尾部能级的技巧
[另见 \citet{Rogers2017}]。随后，\citet{Liu2019} 使用 WW3 对 {\code ST6}
进行了全面重新率定，新的一组参数（即 $a_0,\ a_1,\ a_2,\ B_1$）总结在
表~\ref{tab:ST601} 最后一列。更新后的 {\code ST6} 包不仅能很好地预测常用总体
波浪参数（如有效波高和波周期），还明显改善了高频能级（以饱和谱和均方斜率表示）的估计。
在时长受限测试中，重新率定的 {\code ST6} 得到的全方向频率谱 $E(f)$ 表现出清楚的
过渡行为，即从近似 $f^{-4}$ 幂律过渡到约 $f^{-5}$ 幂律，这与先前现场研究相当
\citep{Forristall1981}。

\begin{table}[htbp]
	\footnotesize
	\begin{center}
	\begin{tabular}{|l|c|c|c|c|} \hline \hline
            {\code ST6} & $a_0$ & $a_1$ & $a_2$ & $B_1$\\
                        & 0.05 & $4.75 \times 10^{-6}$ & $7.00 \times 10^{-5}$ & $6.00 \times 10^{-3}$  \\
            \hline
            {\code GMD} & $\lambda$ & $\mu$ & $\theta_{12}\ (^{\circ})$ & $C_{deep}$  \\
                        & $0.127$ & $0.000$ & $\ 3.0$ & $4.88 \times 10^7$            \\
		        & $0.127$ & $0.097$ & $21.0$ & $1.26 \times 10^8$             \\
		        & $0.233$ & $0.098$ & $26.5$ & $6.20 \times 10^7$             \\
		        & $0.283$ & $0.237$ & $24.7$ & $2.83 \times 10^7$             \\
		        & $0.355$ & $0.183$ & $-$ & $1.17 \times 10^7$                \\
            \hline \hline
	\end{tabular}
	\end{center}
        \caption{{\code ST6} 同 {\code GMD} 非线性求解器 \citep[或具体为 {\code G35};][]{Liu2019} 联用时的率定参数汇总}
	\label{tab:ST602}
	\botline
\end{table}

除了把 {\code ST6} 同 {\code DIA}（第~\ref{sec:NL1} 节）非线性求解器联用外，
\citet{Liu2019} 还尝试将 {\code ST6} 同 $S_{nl}$ 的 {\code GMD} 参数化
（第~\ref{sec:NL3} 节）联用。作为第一步，只采用含 5 个四波组、并采用三参数四波组定义的
{\code GMD} 配置 \citep[即][]{tol:OMOD13d} 中的 {\code G35}。通过
\citet{tol:OMOD13e} 设计的整体遗传优化技术，专门为 {\code ST6} 优化了 {\code GMD}
的可调参数；对应的 {\code ST6} 可调参数总结于表~\ref{tab:ST602}\footnote{
对于表~\ref{tab:ST602} 所示 {\code GMD} 的第五个四波组布局，三参数
$(\lambda,\ \mu,\ \theta_{12})$ 四波组定义退化为双参数 $(\lambda,\ \mu)$ 形式，
而 $\theta_{12}$ 由 $\mu$ 的取值隐含给出 \citep{tol:OMOD13d}。相应地，
{\code ST6+GMD} 与 \citet{tol:OMOD11} 的保守非线性高频滤波器
[另见第~\ref{sec:NLS} 节] 的组合可能是稳定模型积分所必需的，特别是在高分辨率谱网格
（如 $\Delta \theta < 5^{\circ}$）下。},\footnote{不同于表~\ref{tab:ST601}，
这里我们只列出 {\code ST6} 和 {\code GMD} 的重要参数。关于这一特定模型配置中
namelist 变量的详细设置，请读者参阅 {\code bin/README.UoM}。}。
\citet{Liu2019} 表明，在时长受限测试中，{\code GMD} 模拟的 $E(f)$ 与 $S_{nl}$
的精确解（如 {\code WRT}；第~\ref{sec:NL2} 节）高度一致 \citep[另见][]{tol:OMOD13d}。
在 1 年全球后报中，基于 {\code DIA} 的模型会把低频波能量（波周期 $T > 16$ s）
高估 90\%。使用 {\code GMD} 后，此类模型误差显著降低到 $\sim$20\%。值得注意的是，
表~\ref{tab:ST602} 所示 {\code GMD} 方法的计算开销约为 {\code DIA} 的 5 倍。

为灵活控制预报频率区域的高频范围，我们在版本 6.07 中引入
\begin{equation}
f_{hf} = N_{hf}/T_{m0, -1}
\end{equation}
其中 $f_{hf}$ 是截止高频上限 (\ref{eq:tail_E_f})，$T_{m0,-1}$ 是
(\ref{eq:Tm0m1}) 中定义的平均波周期。$N_{hf}$ 通过 namelist 参数 SINFC
设置（表~\ref{tab:ST601}），负的 $N_{hf}$（例如 $N_{hf} = -1$）表示
\emph{高频谱尾在没有任何预设斜率的情况下自由演化。}

\noindent
\paragraph{主导波破碎概率}
按照 \citet[][their Fig.~12]{art:BBY01}，主导波破碎概率 $b_T$ 可由波谱
$F(f, \theta)$ 按如下参数形式估计：
\begin{equation}
b_T = 85.1 \Big[ (\epsilon_p - 0.055) (1 + H_s / d) \Big]^{2.33},
\label{eq:bt}
\end{equation}
其中 $H_s$ 为有效波高，$d$ 为水深，$\epsilon_p$ 为谱峰显著陡度，定义为
\begin{equation}
%\left \lbrace
\begin{array}{rcl}
\epsilon_p &=& H_p k_p / 2\\
H_p &=& 4 \Big[\int_{0}^{2\pi} \int_{0.7f_p}^{1.3f_p} F(f, \theta) \mathrm{d}f \mathrm{d}\theta \Big]^{1/2}
\label{eq:steepness}
\end{array}
%\right.
\end{equation}
其中 $k_p$ 和 $f_p$ 分别为谱峰波数和频率。注意，当我们从 $F(f,\theta)$ 计算
$H_s$、$H_p$ 和 $f_p$（$k_p$）时（式 \ref{eq:bt} 和 \ref{eq:steepness}），
只考虑风海贡献。按照 \citet{Janssen1989} 和 \citet{Bidlot2001}，当
\begin{equation}
\frac{c(k)}{U_{10} \cos (\theta - \theta_u)}  < \beta_w,
\label{eq:beta}
\end{equation}
时，我们把谱分量视为风海；其中 $c(k)$ 为相速，$U_{10}$ 为海面以上 10 m 风速，
$\theta_u$ 为风向，$\beta_w$ 为常数风强迫参数。我们把 $\beta_w$ 实现为可调参数
（{\F MISC} namelist 参数 {\code BTBET}），默认使用 $\beta_w=1.2$。

\noindent
\paragraph{整体调整。} 源项 {\code ST6} 已使用通量参数化 {\code FLX4} 进行了率定。
可通过 {\F FLX4} namelist 参数 {\code CDFAC=1.0E-4}\footnote{自 vers.\,6.07 起，
该值改为 {\code CDFAC=1.0}，因为量级 $10^{-4}$ 已硬编码在 {\code FLX4} 模块中
（表~\ref{tab:ST601}）。} 对阻力参数化 {\code FLX4} 重新缩放，从而实现对风场的
整体调整。其效果类似于调节 {\code ST4} 源项包中变量 $\beta_{max}$，见方程
(\ref{eq:SinWAM4}) 和 (\ref{eq:tauhfint})；该变量可通过 namelist 参数
{\code BETAMAX} 定制（见第 \ref{sec:ST3}--\ref{sec:ST4} 节）。
\citet{pro:Aea11} 和 \citet{art:RA13} 列出了不同取值组合，使模型可适配不同风场。
在优化波浪模型时，建议只重新调节参数 $a_0$、$b_1$ 和 $\mathrm{FAC}$。同样，
$\mathrm{FAC}$ 可能消除风场偏差，而这一偏差通常会随再分析产品选择而变化
\citep[][见表~\ref{tab:ST601}]{Liu2021}。这种削减已在飓风等极端风条件下测试
\citep{art:ZBRY15}。在全球后报中，负输入系数可用于在散点比较中调节波高整体偏差，
而涌浪耗散尺度系数主要影响大海况。当使用离散相互作用近似 ({\code DIA}) 计算四波
相互作用时，比例常数默认值变为 $C=3.00\times10^7$。

\textrm{\textit{\underline{代码限制：}}} 对于版本 4.18 和 5.16，当动力源项积分的最小
时间步远小于总体时间步（即小于 1/15）时，模型会变得不稳定。该问题自 vers.\,6.07 起
已经解决。
